\appendix

\section{Paper--Lean Cross-Reference}

The shared manuscript ownership, notation, and issue register is maintained
in \lean{SEMANTIC-PARITY.md}. The companion coding paper's detailed object
and theorem ledger is \lean{papers/proximity-prize/SEMANTIC_PARITY.md}.
The appendices below record the HCP paper's own declaration and file coverage.

\begin{leanledger}

\subsection{Definitions and Interfaces}

\begin{description}[leftmargin=2.4em,style=nextline]
  \item[Matrix-level public projection and fibers.]
    Sections 2--3; \lean{PublicProjection},
    \lean{HiddenOperatorFiber}, and
    \lean{HiddenOperatorClass} in \lean{HCP.lean}.
  \item[Matrix harmonicity and Hodge instances.]
    Section 2; \lean{Hodge_Instance} and
    \lean{MatrixIsHarmonic} in \lean{HCP.lean}.
  \item[Kernel ambiguity and lossy public data.]
    Sections 2--3; \lean{SameKernel},
    \lean{KernelAmbiguous},
    \lean{PublicCompatibleHodgeAmbiguity}, and
    \lean{LossyPublicProjection} in \lean{HCP.lean}.
  \item[Hodge cipher and operative witness.]
    Section 2; \lean{Hcrypt.HCP.hodgeDecompositionEquiv},
    \lean{Hcrypt.HCP.encrypt}, \lean{Hcrypt.HCP.decrypt},
    \lean{Hcrypt.HCPBridge.Observation},
    \lean{Hcrypt.HCPBridge.HiddenHodgeWitness},
    \lean{Hcrypt.HCPBridge.operativeTarget}, and
    \lean{Hcrypt.HCPBridge.IsOperativeSolution}.
  \item[Representative-indexed public HCP solvers.]
    Section 3; \lean{Hcrypt.HCPBridge.Observation},
    \lean{Hcrypt.HCPBridge.Solver},
    \lean{Hcrypt.HCPBridge.SolverSoundOn}, and
    \lean{Hcrypt.HCPBridge.SolverExactOnHarmonic}. The direct target contract
    is \lean{Hcrypt.HCPBridge.SolverExactOnOperativeTarget}.
  \item[Explicit-lattice CVP comparison.]
    Section 3; \lean{CVP.Instance}, \lean{CVP.Solution}, and
    \lean{CVP.OracleCorrect}.
  \item[Harmonic-lattice bridge.]
    Section 3; \lean{HarmonicCVPBridge} and
    \lean{HarmonicCVPBridge.projectedTarget}.
  \item[Closest-harmonic solution predicates.]
    Section 3; \lean{ClosestHarmonicSolution} and
    \lean{ProjectionClosestHarmonicSolution}.
  \item[Hodge projection model.]
    Sections 2--3; \lean{Laplacian.harmonicProjection},
    \lean{Laplacian.residual}, and \lean{Laplacian.IsHarmonic}.
  \item[Categorical fibration copy surface.]
    Section 4; \lean{IsCartesianMorphism},
    \lean{GrothendieckFibration},
    \lean{CategoricallyEquivalentFibrations}, and
    \lean{FibrationCopy}.
  \item[Typed Helfer/HCP bridge.]
    Sections 2 and 4; \lean{Hcrypt.Helfer.Signature},
    \lean{Hcrypt.Helfer.Interpretation},
    \lean{Hcrypt.Helfer.HomotopyHomomorphism},
    \lean{Hcrypt.Helfer.HomotopyEquivalent},
    \lean{Hcrypt.Helfer.ClosedFormulaInvariance}, and
    \lean{Hcrypt.HCPBridge.Realization}.
  \item[Semantic/public equality surface.]
    Section 4; \lean{Axioms.HelferFOL},
    \lean{Axioms.PublicClassifierFactorization},
    \lean{Axioms.PublicClassifierFactorization.publicFiber_eq_iUnion_baseFibers},
    \lean{Axioms.HomotopyClass}, and \lean{Axioms.Fibration}.
  \item[Heyting/oracle public judgment.]
    Section 4; \lean{Axioms.Heyting3}, \lean{Axioms.OracleQuery},
    \lean{Axioms.OracleAnswer}, and \lean{Axioms.OracleEta2}.
  \item[Concrete finite cellular sheaf and degree-one Hodge operator.]
    Sections 2 and 4; \lean{Sheaf.OrientedTwoComplex},
    \lean{Sheaf.CellularSheaf},
    \lean{Sheaf.CellularSheaf.delta0},
    \lean{Sheaf.CellularSheaf.delta1}, and
    \lean{Sheaf.CellularSheaf.degreeOneLaplacian}.
  \item[Isometric cellular-cochain transport.]
    Section 2; \lean{Hcrypt.Sheaf.CellularSheaf.IsometricCochainEquiv},
    \lean{Hcrypt.Sheaf.CellularSheaf.IsometricCochainEquiv.degreeOneLaplacian_naturality},
    and
    \lean{Hcrypt.Sheaf.CellularSheaf.IsometricCochainEquiv.map_ker_degreeOneLaplacian}.
  \item[Snapshot-bound exact one-cochain.]
    Sections 4--5; \lean{Sheaf.CellularSheaf.SheafStateSnapshot},
    \lean{Sheaf.CellularSheaf.CochainPayloadInterpretation}, and
    \lean{Sheaf.CellularSheaf.AcceptedCommittedOneCochain}.
  \item[Integral sheaf and harmonic message lattice.]
    Sections 2--3; \lean{Sheaf.CellularSheaf.StalkLattice},
    \lean{Sheaf.CellularSheaf.IntegralStructure}, and
    \lean{Sheaf.CellularSheaf.IntegralStructure.harmonicOneLattice}.
    The restriction-preserving integral structure is supplied explicitly.
  \item[Structured Hodge sample and presented harmonic sublattice.]
    Sections 2--3; \lean{Sheaf.CellularSheaf.HodgeNoise},
    \lean{Sheaf.CellularSheaf.HodgeEncryptedCochain}, and
    \lean{Sheaf.CellularSheaf.HarmonicLatticePresentation}.
    A presented full-span sublattice is additional data, not an automatic
    identification of a real harmonic subspace with a discrete lattice.
  \item[Cochain commitment payload.]
    Section 5; \lean{Axioms.CochainCommitmentData},
    \lean{CochainCommitmentData.WellFormed}, and
    \lean{CochainCommitmentData.AddCompatible}.
  \item[Commitment security interface.]
    Section 5; \lean{Axioms.CochainCommitmentAxioms}.
  \item[Proof-carrying public artifact.]
    Section 5; \lean{TCB.HashedCommittedCochain} and
    \lean{TCB.ProofCarryingCommittedCochain}.
  \item[Bound full-sample/exact-payload observation.]
    Sections 2 and 5; \lean{TCB.HCPStackBinding},
    \lean{TCB.BoundHCPObservation}, and
    \lean{TCB.BoundHCPObservation.toObservation}.
    This witness-side record retains the message/opening and a supplied
    context-indexed binding relation; it is not itself a public wire format.
  \item[Projection chain.]
    Section 6; \lean{ProjectionChain.Chain} and
    \lean{ProjectionChain.Chain.toEvaluationProjection}.
  \item[Proximity-stage certificates.]
    Section 6; \lean{Proximity.Gates.SourceMetricGate},
    \lean{Proximity.Gates.HardGateCertificate},
    \lean{Proximity.Gates.ProjectionNaturalityGate}, and
    \lean{Proximity.Gates.NaturalityGateCertificate}. The separately defined
    \lean{Proximity.Gates.SyndromeAdequacyGate} is optional and is not a hard
    gate of the accepted-commitment projection chain.
  \item[Canonical RS/MCA endpoint.]
    Section 6; \lean{CanonicalRS.BridgeData}, \lean{CanonicalRS.MCA.Line},
    and \lean{CanonicalRS.MCA.badChallengeFraction}.
\end{description}

\subsection{Theorems and Status}

\begin{description}[leftmargin=2.4em,style=nextline]
  \item[Kernel ambiguity implies lossy public projection.]
    \lean{kernelAmbiguous_implies_lossy};
    Lean-proved.
  \item[Public-compatible Hodge ambiguity gives kernel ambiguity.]
    \lean{kernelAmbiguous_of_publicCompatibleHodgeAmbiguity};
    Lean-proved.
  \item[Helfer equivalence reaches the HCP public classifier.]
    \lean{Hcrypt.HCPBridge.samePublicClass_of_homotopyEquivalent} and
    \lean{Hcrypt.HCPBridge.publiclyEquivalent_laplacians_of_homotopyEquivalent};
    Lean-proved from the explicit closed-formula invariance and classifier
    factorization hypotheses.
  \item[Typed Hodge-extension separation gives kernel ambiguity.]
    \lean{Hcrypt.HCPBridge.publicCompatibleHodgeAmbiguity_of_kernelSeparation},
    \lean{Hcrypt.HCPBridge.kernelAmbiguous_of_kernelSeparation}, and
    \lean{Hcrypt.HCPBridge.lossyPublicProjection_of_kernelSeparation};
    Lean-proved from an explicit HCP extension-separation witness.
  \item[Hodge encryption and operative target.]
    \lean{Hcrypt.HCP.split_encrypt},
    \lean{Hcrypt.HCP.encrypt_split},
    \lean{Hcrypt.HCP.decrypt_encrypt},
    \lean{Hcrypt.HCP.decrypt_coe_eq_harmonicProjection}, and
    \lean{Hcrypt.HCPBridge.operativeTarget_is_harmonic};
    Lean-proved. The bundled bijection is
    \(\mathcal H\times\mathcal H^\perp\simeq C^1\); the paper separately
    proves the exact/coexact refinement.
  \item[Isometric transport and fixed-coordinate separation.]
    \lean{Hcrypt.Sheaf.CellularSheaf.IsometricCochainEquiv.degreeOneLaplacian_naturality},
    \lean{Hcrypt.Sheaf.CellularSheaf.IsometricCochainEquiv.isHarmonicOneCochain_iff},
    \lean{Hcrypt.HCPBridge.IsometricPresentationEquivalence.harmonic_iff_transport},
    \lean{Hcrypt.HCPBridge.IsometricPresentationEquivalence.map_matrixHarmonicSubspace},
    \lean{Hcrypt.HCPBridge.IsometricPresentationEquivalence.operativeHarmonicComponent_naturality},
    and
    \lean{Hcrypt.HCPBridge.kernelSeparationWitness_of_isometricNonstabilization};
    Lean-proved at the concrete sheaf and abstract HCP interfaces.
  \item[Concrete Hodge ambiguity exists in nontrivial dimension.]
    \lean{Hodge_Ambiguity},
    \lean{constantPublicProjection_publicCompatibleHodgeAmbiguity},
    \lean{constantPublicProjection_kernelAmbiguous}, and
    \lean{constantPublicProjection_lossy};
    Lean-proved.
  \item[Every field-linear image can be a linear kernel.]
    \lean{exists_matrix_ker_eq_range};
    Lean-proved linear-algebra construction.
  \item[Kernel realization for presented generated subspaces.]
    \lean{Hodge_GeneratedSubspace_Reduction} and
    \lean{Hodge_GeneratedSubspace_Reduction_Proof};
    Lean-proved surface.
  \item[Hodge decomposition into harmonic, exact, and coexact parts.]
    \lean{Hodge_Decomposition_Theorem};
    Lean-proved existence and pairwise orthogonality. The paper derives
    uniqueness from those orthogonal summands.
  \item[Ambiguous input contradicts public-data-only Hodge solving.]
    \lean{public_solver_eq_input_of_harmonic} and
    \lean{no_public_solver_on_ambiguous_input};
    Lean-proved lemmas.
  \item[Known-lattice CVP does not select the hidden HCP representative.]
    \lean{Hcrypt.HCPBridge.operativeTargets_ne_of_kernelSeparation},
    \lean{Hcrypt.HCPBridge.no_solver_exactOnOperativeTarget_of_kernelSeparation},
    and \lean{Hcrypt.HCPBridge.cvpOracle_insufficient_for_operativeTarget};
    Lean-proved for the direct operative-target contract. The soundness plus
    harmonic-fixed-point formulation is
    \lean{Hcrypt.HCPBridge.no_solver_on_kernelSeparation} and
    \lean{Hcrypt.HCPBridge.cvpOracle_insufficient_of_kernelSeparation}.
    The older matrix-wide theorems
    \lean{no_uniform_public_hodge_solver} and
    \lean{CVP_Oracle_Insufficient_For_Public_Hodge_Solving} remain
    generic sanity results, not the concrete HCP bridge.
  \item[Projection-model closest-harmonic surface.]
    \lean{matrix_harmonic_iff_projection_isHarmonic},
    \lean{Hodge_Projection_Is_ClosestHarmonic}, and
    \lean{harmonicProjection_is_closestHarmonic};
    Lean-proved.
  \item[Discrete CVP minimizers are preserved by harmonic projection.]
    \lean{HarmonicCVPBridge.squaredDistance_projection_decomposition},
    \lean{HarmonicCVPBridge.isClosestLatticeVector_iff_projectedTarget},
    and \lean{HarmonicCVPBridge.solution_iff_projectedTarget};
    Lean-proved.
  \item[Fibration copies preserve the categorical fibration surface.]
    \lean{copy_categoricallyEquivalent_thm} and
    \lean{FibrationCopy_is_Grothendieck};
    Lean-proved categorical interface.
  \item[Exact sheaf gates supply \(d^{k+1}d^k=0\) and Laplacian-from-coboundary.]
    \lean{Gates.structural_preconditions_of_sheaf_exact_gates};
    Lean-proved alias/interface.
  \item[Concrete cellular coboundaries form a complex.]
    \lean{Sheaf.CellularSheaf.delta1_comp_delta0} and
    \lean{Sheaf.CellularSheaf.isOneCocycle_of_isExactOneCochain};
    Lean-proved from the oriented cell incidences and restriction-composition
    law.
  \item[Concrete degree-one Hodge characterization and orthogonality.]
    \lean{Sheaf.CellularSheaf.degreeOneLaplacian_energy},
    \lean{Sheaf.CellularSheaf.ker_degreeOneLaplacian},
    \lean{Sheaf.CellularSheaf.exactOneSubspace_isOrtho_harmonicOneSubspace},
    \lean{Sheaf.CellularSheaf.harmonicOneSubspace_isOrtho_coexactOneSubspace},
    and \lean{Sheaf.CellularSheaf.exactOneSubspace_isOrtho_coexactOneSubspace};
    Lean-proved.  The separately defined
    \lean{Sheaf.CellularSheaf.kernelNormalizedDegreeOneLaplacian} is proved to
    have the same kernel, not asserted to equal the true Hodge operator or to
    have its spectrum.
  \item[Exact payloads determine openings modulo \(\ker\delta_0\).]
    \lean{Sheaf.CellularSheaf.openingEquivalent_iff_sub_mem_ker},
    \lean{Sheaf.CellularSheaf.decoded_isExactOneCochain}, and
    \lean{Sheaf.CellularSheaf.decoded_isOneCocycle};
    Lean-proved for the concrete cellular-sheaf commitment object.
  \item[Cellular coboundaries preserve the supplied integral structure.]
    \lean{Sheaf.CellularSheaf.IntegralStructure.delta0_mem_integralOneCochains}
    and
    \lean{Sheaf.CellularSheaf.IntegralStructure.delta1_mem_integralTwoCochains};
    Lean-proved from the stalk-lattice restriction-preservation fields.
  \item[Structured samples recover their integral harmonic message.]
    \lean{Sheaf.CellularSheaf.HodgeNoise.value_mem_harmonic_orthogonal} and
    \lean{Sheaf.CellularSheaf.HodgeEncryptedCochain.harmonicOneProjection_observed};
    Lean-proved for the supplied integral sheaf and exact/coexact opening.
    The ambient observed sample is not asserted to be integral or exact.
  \item[Presented harmonic-lattice samples have the encoded point as unique CVP solution.]
    \lean{Sheaf.CellularSheaf.HarmonicLatticePresentation.squaredDistance_observed_decomposition},
    \lean{Sheaf.CellularSheaf.HarmonicLatticePresentation.solution_cvpInstance_iff_eq_point},
    and
    \lean{Sheaf.CellularSheaf.HarmonicLatticePresentation.existsUnique_solution_cvpInstance};
    conditional on an explicit harmonic-lattice presentation and a sample
    constructed from one of its lattice points. This does not supply a
    general CVP algorithm or derive the presentation from public data.
  \item[An actual integral family instantiates the semantic HCP ambiguity.]
    \lean{HCP.AxisFamily.shared_observation_distinct_nonzero_targets},
    \lean{HCP.AxisFamily.no_exact_public_solver},
    \lean{HCP.AxisFamily.cvp_oracle_insufficient}, and
    \lean{HCP.AxisFamily.cvp_solution_iff_message};
    the two cellular sheaves, integral structures, actual Hodge matrices,
    common structured-noise observation and harmonic lattices are constructed.
    Homological-test invariance is proved in the actual bounded-below real
    cochain model. The semantics is a stated fragment, and the oracle
    corollary retains its oracle-existence premise. The observation is
    \((p,b)\), without a full storage-view or security-game claim.
    \lean{weighted_public_observations_eq} and
    \lean{no_exact_encoded_message_decoder} in the same namespace additionally
    separate the explicit integer coefficients \(2\) and \(1\) at one
    shared observation, rather than only separate ambient coordinates.
  \item[The actual bound observation agrees with the semantic operative target.]
    \lean{TCB.HCPRealization.BoundObservationAlignment.operativeTarget_eq_message}
    and \lean{cvp_solution_iff_operativeTarget} in the same namespace;
    derived from actual-operator action and classifier alignment, retaining
    the original sample, payload, snapshot and artifact. A selected richer
    public view requires its own observation-factorization proof.
  \item[The bound observation retains payload/snapshot identity and harmonic recovery.]
    \lean{TCB.BoundHCPObservation.artifact_wellFormed},
    \lean{TCB.BoundHCPObservation.payload_boundToSnapshot}, and
    \lean{TCB.BoundHCPObservation.harmonicProjection_cochainObject};
    Lean-proved from the record's accepted exact component, artifact binding,
    decoded-component equality, and supplied observation-binding relation.
    No modular serializer or concrete public semantic realization follows
    from constructing this conditional record.
  \item[Well-formed payloads have nonempty edge data and supported modulus.]
    \lean{CochainCommitmentData.edge_values_nonempty_of_wellFormed} and
    \lean{CochainCommitmentData.modulus_gt_one_of_wellFormed};
    Lean-proved.
  \item[Proof-carrying artifacts expose the committed payload.]
    \lean{TCB.proofCarrying_cochain_wellFormed};
    TCB boundary theorem.
  \item[Projection chain composes to one evaluation projection.]
    \lean{ProjectionChain.Chain.toEvaluationProjection};
    Lean-proved plumbing.
  \item[Syndrome adequacy turns zero local syndrome into codeword/global candidate.]
    \lean{Proximity.Gates.SyndromeAdequacyGate.codeword_of_wellFormed_zeroSyndrome};
    stage-local interface.
  \item[Stage naturality composes without rechecking.]
    \lean{Proximity.Gates.ProjectionNaturalityGate.of_stageCertificate};
    Lean-proved composition of the cochain/plane and plane/RS restriction and
    syndrome squares.
  \item[Final proximity evidence yields calibrated end-to-end readiness.]
    \lean{Proximity.Gates.CalibratedFinalCertificate.toCalibratedEndToEndReady};
    Lean-proved plumbing.
  \item[Canonical RS list/MCA bounds enter at a cited boundary.]
    \lean{CanonicalRS.CitedBounds.bounds};
    cited external theorem package.
  \item[Canonical MCA bad-challenge fraction is bounded after projection.]
    \lean{CanonicalRS.MCA.projectionBadChallengeFraction_le};
    Lean-proved from cited bounds.
\end{description}

\end{leanledger}

\section{Lean Source Coverage by File}

This appendix is the file-by-file coverage audit for the translation. The body
of the paper gives the mathematical definitions and proofs; this appendix
records where each Lean declaration family enters the written account.
Import-only aggregate files are listed with the files they collect.

\begin{leanledger}

\subsection{Build and Aggregate Files}

\begin{description}[leftmargin=2.4em,style=nextline]
  \item[\lean{lakefile.lean}.]
    Build configuration; no mathematical declaration is translated as a
    theorem.
  \item[\lean{Sheaf.lean}, \lean{Axioms.lean}, \lean{Gates.lean}, \lean{TCB.lean},
    \lean{Consensus.lean}, \lean{Proximity.lean}, and
    \lean{Proximity/CanonicalRS.lean}.]
    Aggregate import surfaces. Their mathematical content is exactly the
    union of the component files listed below.
\end{description}

\subsection{HCP, CVP, and Fibrational Algebra}

\begin{description}[leftmargin=2.4em,style=nextline]
  \item[\lean{HCP.lean}.]
    Section 2 contains the matrix-level HCP vocabulary, public fibers,
    ambiguity hierarchy, typed Helfer realization, representative-bound
    observations and witnesses, the Hodge-coordinate cipher, and the unique
    operative recovery target. It also records isometric presentation
    transport, harmonic-projection naturality, fixed-coordinate
    nonstabilization, classifier nonfaithfulness, and the resulting
    no-exact-representative-decoder theorem. The direct public recovery
    contract is
    \lean{Hcrypt.HCPBridge.SolverExactOnOperativeTarget}.
  \item[\lean{Hardness.lean}.]
    Sections 2--4 contain consequences and comparisons for the HCP interface
    imported from \lean{HCP.lean}. The concrete ambiguity theorems
    \lean{Hodge_Ambiguity}, \lean{Hodge_Ambiguity_Proof},
    \lean{constantPublicProjection_publicCompatibleHodgeAmbiguity},
    \lean{constantPublicProjection_kernelAmbiguous}, and
    \lean{constantPublicProjection_lossy} supply the matrix sanity example
    discussed after \cref{prop:cvp-insufficient}, rather than a concrete typed
    HCP separation witness. The presented-linear-image comparison
    declarations \lean{Hodge_Solver_Correct},
    \lean{exists_matrix_ker_eq_range},
    \lean{Hodge_GeneratedSubspace_Reduction}, and
    \lean{Hodge_GeneratedSubspace_Reduction_Proof} are
    \cref{prop:hcp-to-cvp}. The decomposition theorem
    \lean{Hodge_Decomposition_Theorem} is
    \cref{prop:hodge-decomposition}. The public-solver predicates
    \lean{PublicHodgeInstance}, \lean{PublicHodgeSolver},
    \lean{Hcrypt.HCPBridge.operativeTargets_ne_of_kernelSeparation},
    \lean{Hcrypt.HCPBridge.no_solver_exactOnOperativeTarget_of_kernelSeparation},
    \lean{Hcrypt.HCPBridge.cvpOracle_insufficient_for_operativeTarget},
    \lean{Hcrypt.HCPBridge.no_solver_on_kernelSeparation}, and
    \lean{Hcrypt.HCPBridge.cvpOracle_insufficient_of_kernelSeparation} are
    \cref{prop:cvp-insufficient}. The closest-harmonic predicates
    \lean{ClosestHarmonicSolution}, \lean{ProjectionClosestHarmonicSolution},
    \lean{matrix_harmonic_iff_projection_isHarmonic},
    \lean{Hodge_Projection_Is_ClosestHarmonic}, and
    \lean{harmonicProjection_is_closestHarmonic} record the continuous
    projection surface. The discrete bridge declarations
    \lean{HarmonicCVPBridge},
    \lean{HarmonicCVPBridge.squaredDistance_projection_decomposition},
    \lean{HarmonicCVPBridge.isClosestLatticeVector_iff_projectedTarget}, and
    \lean{HarmonicCVPBridge.solution_iff_projectedTarget} are
    \cref{prop:hodge-projection-is-cvp}. The categorical declarations
    \lean{IsCartesianMorphism}, \lean{GrothendieckFibration},
    \lean{CategoricallyEquivalentFibrations}, \lean{Wrapper}, \lean{Fibration},
    \lean{FibrationCopy}, \lean{ToWrapper}, \lean{FromWrapper},
    \lean{copy_categoricallyEquivalent_thm}, and \lean{FibrationCopy_is_Grothendieck}
    are \cref{prop:fibration-copy-equivalence}.
  \item[\lean{CVP.lean}.]
    Section 3 translates \lean{AmbientVector}, \lean{ExplicitLattice},
    \lean{latticeVector},
    \lean{InExplicitLattice}, \lean{IsClosestLatticeVector},
    \lean{squaredEuclideanDistance}, \lean{Instance}, \lean{Solution},
    \lean{Oracle}, \lean{LatticeOracleCorrect}, \lean{OracleCorrect},
    \lean{SolverCorrect}, \lean{solverCorrect_iff_oracleCorrect},
    \lean{HardnessAssumption}, and
    \lean{zero_lattice_oracle_correct} as the known-lattice comparison
    surface.
  \item[\lean{Helfer/Signature.lean},
    \lean{Helfer/HomotopyEquivalence.lean},
    \lean{Helfer/Invariance.lean}, \lean{Axioms/HelferFOL.lean},
    \lean{Axioms/Fibration.lean}, and
    \lean{Axioms/HomotopyClass.lean}.]
    Section 4 formalizes multi-sorted signatures, finite-product
    interpretations, coherent homotopy homomorphisms, and componentwise
    homotopy equivalence for cofibrant-fibrant interpretations. It also
    extracts the Helfer-derived path/diagonal and fibrational interface, then
    adds the project-specific public classifiers and public fibers, classifier
    factorization through a base projection, class-membership witnesses,
    transport, semantic well-formedness clauses, fiber systems, ephemeral
    fiber specifications, and nonfaithful classifiers. This is not a complete
    formalization of Helfer's model-categorical construction or
    formula-in-context invariance proof; the closed-formula consequence of
    Theorem~12.3 is an explicit Lean hypothesis. The local equivalence, symmetry, transitivity,
    membership, and nonfaithfulness lemmas justify the public fiber language
    used in HCP.
\end{description}

\subsection{Foundational, Sheaf, Spectral, and Judgment Axioms}

\begin{description}[leftmargin=2.4em,style=nextline]
  \item[\lean{Sheaf/CellComplex.lean}, \lean{Sheaf/Cellular.lean},
    \lean{Sheaf/Hodge.lean}, \lean{Sheaf/Equivalence.lean}, and
    \lean{Sheaf/Commitment.lean}.]
    Sections 2, 4, and 5 translate the concrete finite oriented two-complex;
    finite real coordinate stalks and face-to-coface maps; the typed spaces
    \(C^0,C^1,C^2\); concrete \(\delta_0\) and \(\delta_1\); the proof that
    \(\delta_1\delta_0=0\); degree-zero global sections, degree-one cocycles,
    and exact degree-one cochains; the true degree-one Hodge operator
    \(\delta_0\delta_0^*+\delta_1^*\delta_1\); its energy, kernel, and
    orthogonality theorems; the distinct kernel-normalized operator;
    degreewise isometric cochain equivalences and their adjoint, Laplacian,
    harmonic-subspace, and kernel transport theorems; immutable
    sheaf-state snapshots; snapshot-bound payload interpretation; accepted
    exact one-cochains; and the quotient of openings by \(\ker\delta_0\).
  \item[\lean{Sheaf/Integral.lean}.]
    Sections 2--3 and 5 use \lean{StalkLattice}, its additive carrier and
    \lean{Preserves} predicate, and \lean{IntegralStructure}. The latter
    selects integral stalk lattices preserved by the cellular restrictions.
    The induced \lean{integralZeroCochains}, \lean{integralOneCochains}, and
    \lean{integralTwoCochains} are additive subgroups;
    \lean{harmonicOneLattice} is the integral degree-one subgroup intersected
    with the harmonic subspace. The two integral-coboundary preservation
    theorems are proved for this supplied structure. No canonical modular
    encoding/decoding is constructed here.
  \item[\lean{Sheaf/Distance.lean} and \lean{Sheaf/CochainSemantics.lean}.]
    Distance to the actual exact subspace is nonnegative, vanishes exactly
    on exact cochains, and is stable under Euclidean perturbations. The
    semantics adapter intertwines the actual cochain operators and decoder,
    identifies snapshot binding, and derives exactness, cocycle and acceptance
    transport. Other legacy ambient gates remain explicit. The carrier
    equivalence is not asserted to preserve norms.
  \item[\lean{Sheaf/HarmonicLattice.lean}.]
    Sections 2--3 use \lean{HodgeNoise}, \lean{HodgeEncryptedCochain},
    \lean{observed}, and harmonic projection recovery for
    \(b=h+\delta_0a+\delta_1^*c\).
    \lean{HarmonicLatticePresentation} additionally supplies an explicit
    full-span harmonic sublattice, membership proofs and its real-span
    equality. Its \lean{point}, \lean{encrypt}, and \lean{cvpInstance}
    constructions lead to the squared-distance decomposition and unique
    closest-point result for those encoded samples. The presentation need
    not contain every integral harmonic cochain. Noise sampling, a secrecy
    game and modular wire realization remain separate obligations.
  \item[\lean{Sheaf/ObservationQuotient.lean}.]
    The full kernel-normalized residual equals the Hodge-noise value.
    \Cref{prop:exact-component-obstruction} is formalized by the coexact-ray,
    residual-nonfactorization and far-preservation obstruction theorems.
    Their source carrier is unrestricted Hodge noise, and the extraction
    must factor through its decoded exact component; no raw-byte
    factorization or unrestricted production-profile claim is inferred.
  \item[\lean{Axioms/Foundation.lean} and \lean{Gates/Exact.lean}.]
    Section 4 translates \lean{Foundation} and
    \lean{FoundationWellFormed}: the cochain-complex premise, coboundary
    linearity, Hodge--Laplacian formation, Hodge decomposition, and Hodge
    orthogonality.  Commitment authentication/hiding, topology connectivity,
    spectral policy, and key-image/freshness conditions are separate
    interfaces rather than fields of \lean{Foundation}. The exact-gate
    aliases \lean{ExactGateBundle}, \lean{SheafExactGateBundle},
    \lean{FoundationWithSheafExactGates}, and
    \lean{structural_preconditions_of_sheaf_exact_gates} are
    \cref{prop:exact-structural-preconditions}.
  \item[\lean{Axioms/SheafCoherence.lean}.]
    Section 4 translates restriction-shape compatibility, edge naturality,
    triangle-face commutativity, sheaf cochain complexes, Hodge coherence,
    finite simplicial nerves, \lean{SheafCoherence},
    \lean{SheafCoherenceWellFormed},
    typed degree-zero through degree-two coboundaries,
    \lean{CommittedCochainWellFormed},
    \lean{RepresentsExactOneCochain}, \lean{IsOneCocycle},
    \lean{oneCocycle_of_committedCochainWellFormed},
    \lean{FoundationWithSheafCoherence}, and
    \lean{coherence_implies_structural_preconditions}.
  \item[\lean{Axioms/Grassmannian.lean}.]
    Sections 4 and 6 translate the Grassmannian metric and geometry objects:
    \lean{GrassmannianDistance}, \lean{chordal_leq_geodesic},
    \lean{KarcherMean}, \lean{SheafLaplacian}, \lean{CechNerve},
    \lean{PluckerEligibility}, \lean{ConsensusGeometry}, and
    \lean{ConsensusGeometryWellFormed}.
  \item[\lean{Axioms/SpectralBounds.lean} and \lean{Gates/Spectral.lean}.]
    Section 4 translates the constants
    \(\sqrt2/2\) and \(2\sqrt2\), the predicates
    \lean{SpectralGapInPoVBounds}, \lean{NearPoVSpectralBoundary},
    \lean{TsirelsonPredicate}, \lean{NearTsirelsonBoundary}, the aliases
    \lean{SpectralGapGate} and \lean{NearSpectralBoundary}, the geometry
    witness \lean{ByConstructionPoVGeometry}, and the theorems
    \lean{byConstruction_gap_in_bounds} and
    \lean{byConstruction_passes_spectral_gate}.
  \item[\lean{Tsirelson.lean}.]
    Section 4 translates the Grassmannian/Laplacian family
    \lean{Grassmannian_Laplacians}, \leanOmega{_G},
    \lean{LaplacianSpace}, \lean{SpectralGap}, \leanOmega{_1},
    \lean{projector_map}, \lean{Classical_Strategy},
    \lean{Quantum_Strategy}, and \lean{MaxEigenvalue}; the infinitude and
    injectivity results \lean{R_infinite}, \lean{infinite_lines},
    \lean{Grassmannian_Laplacian_injective},
    \lean{projector_injective}, \lean{laplacian_injective},
    \leanOmega{_1_infinite}, \lean{R_uncountable},
    \leanOmega{_1_infinite_v2}, and \leanOmega{_1_Nonempty}; the Laplacian-space
    inclusions \lean{laplacian_properties},
    \lean{laplacian_is_hermitian},
    \leanOmega{_1_subset_LaplacianSpace},
    \leanOmega{_1_subset_LaplacianSpace_v3},
    \leanOmega{_1_subset_LaplacianSpace_v5},
    \leanOmega{_1_subset_LaplacianSpace_explicit_v2}, and
    \leanOmega{_1_subset_LaplacianSpace_explicit}; the classical/quantum
    statement forms \lean{Classical_Bound_Statement} and
    \lean{Quantum_Bound_Statement}; and the empty-family sanity checks
    \lean{Classical_Bound_empty_family} and
    \lean{Quantum_Bound_empty_family}.
  \item[\lean{Axioms/Heyting.lean} and
    \lean{Axioms/OracleModality.lean}.]
    Section 4 translates the three-valued Heyting algebra, meet, join,
    implication, negation, failure of double-negation elimination, abstract
    Heyting collapse, nonfaithful public truth, spectral judgments, consensus
    actions, oracle queries and answers, \(\Omega_{\mathrm{nv}}\), witness
    collapses, quotient equivalence, sheafification, oracle forcing,
    \(\eta_1/\eta_2\) oracle modalities, Lawvere--Tierney idempotence,
    boundary judgment behavior, canonical consensus admissibility, oracle
    pipelines, and the oracle/consensus bridge theorems.
  \Needspace{6\baselineskip}
  \item[\lean{Axioms/StageAdmissibility.lean},
    \lean{Gates/Stage.lean}, and \lean{Gates/Oracle.lean}.]
    This appendix records the stage and oracle aliases: pipeline stages, stage
    admissibility, transaction-critical invariant bundles, Python-stage
    snapshots, refinement predicates, and the theorem that a refined
    transaction-critical pass is admissible.
\end{description}

\subsection{Laplacian Geometry and Search Spaces}

\begin{description}[leftmargin=2.4em,style=nextline]
  \item[\lean{Laplacian/Hodge.lean} and \lean{Laplacian.lean}.]
    Sections 2--3 translate harmonic projection, residual vector, residual
    norm, harmonicity, projection membership, nonnegativity, and all
    equivalences between zero residual, membership in the harmonic subspace,
    and harmonicity.
  \item[\lean{Laplacian/Projection.lean}.]
    Section 3 translates the projection Laplacian
    \(\Delta_W=I-P_W\) and its pointwise application formula.
  \item[\lean{Laplacian/Spectral.lean}.]
    Section 4 translates sheaf Laplacians, spectral fingerprints, spectral
    equivalence, scalar Laplacians, kernel and eigenvalue lemmas, scalar
    fingerprints, the legacy relation defined by equal kernel dimension, and
    the theorem that spectrum is not determined by nullity. That relation
    does not assert Helfer homotopy equivalence.
  \item[\lean{Laplacian/Likelihood.lean}.]
    Section 4 records classical selector and density-based likelihood
    interfaces, Gaussian likelihoods, minimum spectral gaps, integrated
    likelihood scores and Gaussian factorization lemmas. Legacy names such
    as ``quantum adversary'' and \lean{success_prob_*} supply neither a
    quantum computation model nor HCP recovery probabilities. The corrected
    Euclidean integrated scores are not identified with the legacy
    supremum-norm integrals. These interfaces are outside the selected entry.
  \item[\lean{Laplacian/Grassmannian.lean},
    \lean{Laplacian/Grassmannian/Injectivity.lean},
    \lean{Laplacian/Grassmannian/SearchSpace.lean}, and
    \lean{Laplacian/Grassmannian/ParametricFamily.lean}.]
    Sections 4 and 6 translate harmonic subspaces, the map to projection
    Laplacians, the families \(\Omega_r\) and \(\Omega_1\), membership in
    Laplacian space, nonemptiness of \(\Omega_k\), injectivity of the
    Grassmannian Laplacian map, and the parametric construction of families of
    \(k\)-planes.
\end{description}

\subsection{Commitments, Public Artifacts, and TCB Boundary}

\begin{description}[leftmargin=2.4em,style=nextline]
  \item[\lean{Axioms/CochainCommitment/Payload.lean},
    \lean{Axioms/CochainCommitment/Blinded.lean},
    \lean{Axioms/CochainCommitment/Security.lean}, and
    \lean{Axioms/CochainCommitment.lean}.]
    Section 5 translates the canonical modulus, byte lengths, wire length,
    \lean{CochainCommitmentData}, edge-count matching, nonempty edge data,
    supported modulus, well-formedness, commitment sums, common-topology and
    additive compatibility, zero and addition operations, payload lemmas,
    exact-payload authentication, wrapper hiding, context binding,
    additivity/security interfaces, the typed exact-payload to
    Hodge-blinded-value bridge, exact--harmonic orthogonality and decomposition
    theorems, blinded and vector commitment axioms, cohomological MAC axioms,
    and the canonical commitment construction from foundational geometry plus
    separately supplied commitment-specific hypotheses.
  \item[\lean{TCB/ProofCarrying.lean}, \lean{TCB/Runtime.lean},
    \lean{TCB/SNARK.lean}, \lean{TCB/CompressionBoundary.lean},
    \lean{TCB/ObservationBoundary.lean}, and
    \lean{TCB/WitnessVerifier.lean}.]
    Section 5 and this appendix record hash words, hashed committed cochains,
    proof-carrying cochains, runtime specification boundaries, runtime
    acceptance predicates, SNARK-carrying cochains, compression contracts,
    compressed public data, observation boundaries, gap-decision certificates,
    witness-compression envelopes, public state updates, and theorems stating
    that accepted artifacts expose well-formed cochain and oracle predicates.
  \item[\lean{Helfer/CochainModel.lean}, \lean{Sheaf/HCPAxisFamily.lean},
    \lean{HCP/AxisFamily.lean}.]
    Section 3's concrete family. The first module constructs the real cochain
    Quillen model, retains the actual cellular differentials and proves
    homological-test invariance for all typed interpretations. The second
    constructs integral sheaves, samples, operator matrices, chain isometry,
    spectrum and harmonic lattices. The third derives the semantic realization,
    common observation, distinct nonzero targets and both positive and negative
    CVP comparisons from those same objects. No full first-order semantics,
    cryptographic sampling experiment or runtime profile is constructed.
  \item[\lean{TCB/HCPRealization.lean}, \lean{TCB/HCPAxisFamily.lean}.]
    Sections 3 and 5's application join. Actual matrix action gives the
    cellular kernel/projector equality; snapshot-classifier agreement gives
    the original observation's hidden witness and operative message. The
    existing CVP comparison is transported to that target. The concrete family
    instantiates the operator alignment and can retain an original artifact
    under supplied acceptance, decoding and envelope bindings. The actual
    computed positive gap also satisfies the spectral gate. Complete public
    storage-view equality is not inferred from the class equality.
  \item[\lean{TCB/HCPObservation.lean}.]
    Sections 2 and 5 use \lean{HCPStackBinding} and
    \lean{BoundHCPObservation} to join the full Hodge sample to its accepted
    exact committed component, artifact leaf digest, immutable snapshot root
    and lifecycle context. The \lean{latticeObject}, \lean{cochainObject},
    \lean{hashObject}, \lean{snapshotRoot}, and \lean{toObservation} accessors
    retain those same objects. Artifact well-formedness, snapshot binding and
    harmonic recovery are proved from the explicit record fields. The
    classifier, decoder and observation-binding relation are supplied;
    constructing the record does not instantiate their public realization.
    The record contains private witness data and is not a serialization or
    public disclosure prescription.
  \item[\lean{Axioms/ZKP.lean}.]
    Section 5 records the zero-knowledge proof object and its well-formedness
    predicate as an assumed proof-carrying artifact interface.
\end{description}

\subsection{Projection Chain, Proximity Gates, and RS/MCA Endpoint}

The universal standalone coding results are owned by the companion Paper 3
and its detailed source ledger. Its completed generic distance, real-shell,
phase-split and all-parameter family constructions reside in
\lean{Proximity/Core/Distance.lean},
\lean{Proximity/Core/RadiusShell.lean},
\lean{Proximity/Capacity/Realization.lean},
\lean{Proximity/Capacity/ListPhase.lean},
\lean{Proximity/ListDecoding/PhaseSplit.lean},
\lean{Proximity/MCA/RealThreshold.lean}, and
\lean{Proximity/Application/DistanceCalibration.lean}.
This appendix records their HCP-specific consumers without repeating the
standalone coding proofs.

\begin{description}[leftmargin=2.4em,style=nextline]
  \item[\lean{Proximity/Problem.lean}, \lean{Proximity/Public.lean},
    \lean{Proximity/Projection.lean}, \lean{Proximity/PluckerShadow.lean},
    \lean{Proximity/GrassmannianPlucker.lean}, and
    \lean{Proximity/PluckerCochain/PublicBinding.lean}.]
    Section 6 records the problem-facing public projection vocabulary, the
    abstract Pl\"ucker coordinate-system and public-shadow contracts, and the
    interfaces that prepare the concrete Pl\"ucker coordinate and finite-code
    datum.
  \item[\lean{Proximity/Application/GapBoundary.lean}.]
    Section 6.1 formalizes disjointness of the closed promises under a strict
    gap, the equal-radius boundary obstruction, and strict-gap views retaining
    the existing carrier, score and radii. The classical specification solver
    characterizes logical consistency; it supplies neither an executable or
    efficient solver nor a concrete application profile.
  \item[\lean{Proximity/Grassmannian.lean},
    \lean{Proximity/GrassmannianPlucker.lean}, and
    \lean{Proximity/GrassmannianPlucker/Concrete.lean}.]
    Section 6 translates harmonic Grassmannian points, exterior coordinates,
    frame witnesses, decomposability, framed systems, raw and normalized
    Pl\"ucker coordinates, soundness, projective equivalence, projective
    normalization, projective framed systems, and projective sound inputs.
  \item[\lean{Proximity/Concrete/GrassmannianPlucker/Reconstruction.lean}
    and \lean{TCB/PluckerObservation.lean}.]
    Section 6 proves both concrete exterior-algebra soundness inputs,
    reconstructs the operative subspace and harmonic projector from a genuine
    full projective point, and proves full-shadow equality exactly when the
    ranked points agree. Identification with \(\ker\Delta_x\) gives the
    operative harmonic target; digest inversion and efficient real-coordinate
    algorithms are outside these results.
  \item[\lean{Proximity/CochainBridge.lean},
    \lean{Proximity/PluckerCochain.lean},
    \lean{Proximity/PluckerCochain/Concrete.lean}, and
    \lean{Proximity/PluckerCochain/PublicBinding.lean}.]
    Section 6 translates valid frames, Pl\"ucker-to-cochain instances, local
    syndrome targets, codeword predicates, evaluation projections, invariant
    checklists, payload encoders, commitment-sum residuals, payload
    realizations, hashed Pl\"ucker coordinates, observed frame bindings, and the
    theorems that valid Pl\"ucker coordinates bind to sheaf-coherent exact
    one-cochains, hence to accepted committed cocycles and local codewords.
  \item[\lean{Proximity/CochainPlane.lean},
    \lean{Proximity/CochainPlane/Core.lean},
    \lean{Proximity/CochainPlane/Concrete.lean},
    \lean{Proximity/CochainPlane/ReedSolomon.lean}, and
    \lean{Proximity/CochainReedSolomon.lean}.]
    Section 6 translates plane projections, plane extraction contexts,
    local nondegeneracy, alignment, projected planes, plane targets,
    plane-to-evaluation projections, typed restriction and syndrome naturality
    at the cochain/plane and plane/RS transitions,
    local-codeword-to-target-codeword theorems, separate plane and RS metric
    handoffs, faithful-on-region plane encodings, interpretations, and RS
    received words.
  \item[\lean{Proximity/ProjectionChain.lean} and
    \lean{Proximity/ProjectionChain/Core.lean}.]
    Section 6 translates the complete projection chain, including accessors
    for committed cochain, plane projection, received word, and the composite
    evaluation projection.
  \item[\lean{Proximity/Gates.lean}, \lean{Proximity/Gates/Stage.lean},
    \lean{Proximity/Gates/Final.lean}, and
    \lean{Proximity/GateEvidence.lean}.]
    Section 6 translates only proximity gates: source metric, Pl\"ucker metric,
    semantic alignment, optional syndrome adequacy, projective-coordinate
    respect, three-stage noncollapse through cochain and plane radii,
    far-region separation, runtime binding, structural bundles, hard-gate and
    derived naturality certificates, calibrated final certificates, nested
    witness-local state buses, stage checklists, stage certificates, radius
    calibrations, final readiness, evidence slots, evidence databases, and the
    binding/completeness lemmas for proximity evidence.
  \item[\lean{Proximity/RS/ListDecoding.lean},
    \lean{Proximity/RS/ReedSolomon.lean},
    \lean{Proximity/RS/CorrelatedAgreement.lean},
    \lean{Proximity/RS/Target.lean}, and \lean{Proximity/RS.lean}.]
    Section 6 translates words, block codes, Hamming and relative Hamming
    balls, list-size bounds, interleaved and additive codes, RS parameters,
    vector codewords, correlated-agreement and mutual-correlated-agreement
    predicates, bad-parameter sets, bridge conditions, calibrated
    \(\HE\)-to-RS bridges, and pullback-correctness theorems.
  \item[\lean{Proximity/CanonicalRS/Parameters.lean},
    \lean{Proximity/CanonicalRS/PowerDomain.lean},
    \lean{Proximity/CanonicalRS/Encoding.lean},
    \lean{Proximity/CanonicalRS/Metric.lean},
    \lean{Proximity/CanonicalRS/Target.lean},
    \lean{Proximity/CanonicalRS/FarSeparation.lean},
    \lean{Proximity/CanonicalRS/BridgeData.lean},
    \lean{Proximity/CanonicalRS/BridgeWitness.lean},
    \lean{Proximity/CanonicalRS/BridgeCertificate.lean},
    \lean{Proximity/CanonicalRS/CitedBounds.lean},
    \lean{Proximity/CanonicalRS/MCA.lean},
    \lean{Proximity/CanonicalRS/MCA/Line.lean}, and
    \lean{Proximity/CanonicalRS/MCA/Concrete.lean}.]
    Section 6 translates the canonical finite field, length, degree bound,
    rate, Johnson and capacity radii, smooth domains, power-root domain,
    canonical encoding, the proof that the unbounded integer encoding is not
    globally injective, Hamming-to-code metrics, target bounds, far-separation
    witnesses, bridge data and theorem slots, bridge certificates, cited
    BCHKS-style bounds, MCA block code, challenge and bad-challenge fractions,
    and the bound on the projected bad-challenge fraction.
\end{description}

\newpage
\subsection{Subsequent Protocol and Implementation Boundary}

The bounded state/transition, candidate-indexed evidence and typed storage
modules under \lean{Protocol/Ledger/}, \lean{Protocol/Evidence/}, and
\lean{Protocol/Storage/} belong to Paper 2;
their completion status is recorded in \lean{papers/paper2/SCOPE.md}.
They are downstream consumers of the mathematical objects above.

\begin{description}[leftmargin=2.4em,style=nextline]
  \item[\lean{Axioms/Consensus.lean}, \lean{Axioms/Ledger.lean},
    \lean{Axioms/Coregulator.lean}, and
    \lean{Axioms/PythonHomotopyRefinement.lean}.]
    This appendix records consensus, ledger, coregulator, and Python homotopy
    refinement as boundary objects. Their full protocol semantics belong to
    the subsequent protocol paper, but their predicates are named here so the
    Paper 1 mathematical objects have explicit downstream interfaces.
  \item[\lean{Consensus/Geometry.lean},
    \lean{Consensus/Reliability.lean}, and
    \lean{Consensus/ProofOfValue.lean}.]
    This appendix records Pl\"ucker eligibility and consensus geometry aliases,
    reliability evidence and bounded reliability scores, proximity quality,
    Proof-of-Value gates, Proof-of-Value inputs, scores with and without
    reliability/proximity, and the score bounds. These are protocol-facing
    uses of the Section 4 and Section 6 mathematical predicates.
  \item[\lean{Chain.lean}.]
    This appendix records validators, transactions, blocks, protocol states, spectral
    bounds, topology and freshness predicates, block validity, tips, heights,
    extension, transaction and block acceptance, ledger-state updates,
    key-image preservation theorems, validator eligibility, cumulative gap,
    preference, transaction composition and identities, quadratic weight,
    Sybil-resistance weighting, chain-growth monotonicity, fork-resolution
    transitivity, and connectedness of eligible validators. The detailed
    protocol architecture remains in the subsequent protocol paper.
\end{description}

\end{leanledger}
